\documentclass[landscape,11pt]{article}

\usepackage[utf8]{inputenc}
\usepackage[T2A]{fontenc}
\usepackage[russian]{babel}
\usepackage{geometry}
\usepackage{tikz}
\usetikzlibrary{arrows.meta, positioning}

\pagestyle{empty}
% Высоты A4 недостаточно для двух симметричных частей диаграммы.
\geometry{paperwidth=420mm,paperheight=297mm,left=12mm,right=12mm,top=10mm,bottom=10mm}

\tikzset{
	uml/.style={
		draw,
		rounded corners=1pt,
		align=left,
		font=\scriptsize,
		text width=43mm,
		inner sep=3pt,
		minimum height=12mm
	},
	relation/.style={
		-{Stealth[length=2mm]},
		line width=.45pt
	},
	dependency/.style={
		-{Stealth[length=2mm]},
		dashed,
		line width=.45pt
	}
}

\begin{document}
	\centering
	\begin{tikzpicture}[x=1cm, y=0.7cm]
		\node[font=\Large\bfseries] at (-1,6.5)
		{UML-диаграмма классов клиент-серверного приложения};
		
		% ===================== Клиентская часть =====================
		
		% (бывший ClientMain на (-15,4) → опущен на -1 и перемещён на место ClientConsole)
		\node[uml] (ClientMain) at (-15,-1)
		{
			\textbf{client::Main}\\
			\hrulefill\\
			+ main(String[]): void
		};
		
		% (ClientConsole удалён; на его место поставлен ClientRunner)
		\node[uml] (ClientRunner) at (-15,3)
		{
			\textbf{client::ClientRunner}\\
			\hrulefill\\
			- credentials: Credentials\\
			- client: DatagramClient\\
			+ run(), send(Request), authorize()
		};
		
		\node[uml] (ClientCommandManager) at (-1,3)
		{
			\textbf{client::CommandManager}\\
			\hrulefill\\
			- commands: Map<String,\\
			Command>\\
			+ getCommands(): Map
		};
		
		\node[uml] (DatagramClient) at (-8,-1)
		{
			\textbf{client::DatagramClient}\\
			\hrulefill\\
			- channel: DatagramChannel\\
			- server: SocketAddress\\
			+ request(Request): Response\\
			+ close(): void
		};
		
		\node[uml] (ClientCommand) at (-1,-1)
		{
			\textbf{client::commands::}\\
			\textbf{Command}\\
			\hrulefill\\
			+ execute(argument, runner)\\
			\hrulefill\\
			Register, Login, Help, Add,\\
			Update, ExecuteScript, \ldots
		};
		
		% ===================== Общие классы =====================
		
		\node[uml] (DatagramTransfer) at (-8,-7)
		{
			\textbf{common::DatagramTransfer}\\
			\hrulefill\\
			+ send(channel, data, addr): void\\
			+ receive(channel, timeout): ReceivedData\\
			+ receiveAvailable(...): ReceivedData
		};
		
		\node[uml] (request) at (-15,-7)
		{
			\textbf{common::Protocol.Request}\\
			\hrulefill\\
			- command: CommandType\\
			- payload: Payload\\
			- credentials: Credentials
		};
		
		\node[uml] (response) at (-1,-7)
		{
			\textbf{common::Protocol.Response}\\
			\hrulefill\\
			- ok: boolean\\
			- payload: Payload
		};
		
		\node[uml] (Model) at (6,-7)
		{
			\textbf{common::Model.}\\
			\textbf{SpaceMarine}\\
			\hrulefill\\
			- id, name, health\\
			- coordinates: Coordinates\\
			- creationDate: ZonedDateTime\\
			- chapter: Chapter\\
			+ compareTo(SpaceMarine)
		};
		
		% ===================== Серверная часть =====================
		
		\node[uml] (ServerMain) at (-15,-14)
		{
			\textbf{server::Main}\\
			\hrulefill\\
			+ main(String[]): void
		};
		
		\node[uml] (ServerConsole) at (-15,-20)
		{
			\textbf{server::ServerRunner}\\
			\hrulefill\\
			- console: Console\\
			+ run(stopServer): void
		};
		
		\node[uml] (datagram) at (-8,-14)
		{
			\textbf{server::DatagramServer}\\
			\hrulefill\\
			- channel: DatagramChannel\\
			- sendingPool: ExecutorService\\
			+ start(), run(), stop()\\
			+ reader Thread, processor Thread\\
			+ cached pool for responses
		};
		
		\node[uml] (handler) at (-8,-20)
		{
			\textbf{managers::RequestHandler}\\
			\hrulefill\\
			- userRepository:\\
			UserRepository\\
			+ handle(Request): Response
		};
		
		\node[uml] (command) at (-1,-14)
		{
			\textbf{commands::Command}\\
			\hrulefill\\
			+ apply(Request,\\
			AuthenticatedUser)\\
			\hrulefill\\
			Add, Update, RemoveById,\\
			Clear, RemoveLower, \ldots
		};
		
		\node[uml] (service) at (6,-14)
		{
			\textbf{managers::}\\
			\textbf{CollectionManager}\\
			\hrulefill\\
			- byId: ConcurrentMap\\
			- collection: ConcurrentSkipListSet\\
			- mutationLock: ReentrantLock\\
			+ add(), update(), clear()\\
			+ snapshot(), descending()
		};
		
		\node[uml] (storage) at (6,-20)
		{
			\textbf{managers::}\\
			\textbf{DatabaseManager}\\
			\hrulefill\\
			+ openConnection():\\
			Connection\\
			\textbf{MarineRepository}\\
			\textbf{UserRepository}
		};
		
		\node[uml] (manager) at (-1,-20)
		{
			\textbf{managers::}\\
			\textbf{CommandManager}\\
			\hrulefill\\
			- commands: Map\\
			- history: Deque\\
			+ register(), getCommand()
		};
		
		\node[uml] (postgres) at (13,-20)
		{
			\textbf{PostgreSQL}\\
			\hrulefill\\
			lab7\_users\\
			lab7\_space\_marines\\
			справочные таблицы
		};
		
		% ===================== Связи: клиент =====================
		
		\draw[dependency]
		(ClientMain) -- node[left, font=\scriptsize]{создаёт} (ClientRunner);
		
		\draw[dependency]
		(ClientRunner) -- node[right, font=\scriptsize]{использует} (DatagramClient);
		
		\draw[dependency]
		(ClientCommandManager) -- node[right, font=\scriptsize]{хранит} (ClientCommand);
		
		% Новые связи (как у сервера)
		\draw[dependency]
		(ClientMain) -- node[above, font=\scriptsize]{создаёт} (DatagramClient);
		
		\draw[dependency]
		(DatagramClient) -- node[above, font=\scriptsize]{использует} (ClientCommandManager);
		
		% Связь ClientRunner → ClientCommand удалена по требованию
		
		% ===================== Связи: сервер =====================
		
		\draw[dependency]
		(ServerMain) -- node[left, font=\scriptsize]{создаёт} (handler);
		
		\draw[dependency]
		(ServerMain) -- node[above, font=\scriptsize]{запускает} (datagram);
		
		\draw[dependency]
		(ServerMain) -- node[left, font=\scriptsize]{запускает} (ServerConsole);
		
		\draw[dependency]
		(datagram) -- node[right, font=\scriptsize]{использует} (handler);
		
		\draw[dependency]
		(handler) -- node[above, font=\scriptsize]{использует} (manager);
		
		\draw[dependency]
		(manager) -- node[below, font=\scriptsize]{хранит} (command);
		
		\draw[relation]
		(command) -- node[above, font=\scriptsize]{использует} (service);
		
		\draw[dependency]
		(service) -- node[right, font=\scriptsize]{читает/изменяет} (storage);
		
		\draw[relation]
		(storage) -- node[above, font=\scriptsize]{JDBC} (postgres);
		
		% ===================== Связи: UDP-обмен =====================
		
		\draw[dependency]
		(datagram) -- node[right, font=\scriptsize]{использует} (DatagramTransfer);
		
		\draw[dependency]
		(DatagramClient) -- node[right, font=\scriptsize]{использует} (DatagramTransfer);
		
		\draw[dependency]
		(datagram) -- node[right, font=\scriptsize]{десериализует} (request);
		
		\draw[dependency]
		(datagram) -- node[right, font=\scriptsize]{сериализует} (response);
		
		\draw[dependency]
		(DatagramClient) -- node[above, font=\scriptsize]{сериализует} (request);
		
		\draw[dependency]
		(DatagramClient) -- node[right, font=\scriptsize]{десериализует} (response);
		
		
	\end{tikzpicture}
	
\end{document}